\subsection{El predictor de un residual y su identidad diferencial}\label{sec:B3T01}
En las fórmulas de esta sección la letra $B$ designa, además del tablero residual, la matriz de
restricciones de filas y columnas; escribimos RC por «filas y columnas». En cada componente una
ecuación de masa es redundante: al sumar por filas o por columnas se obtiene la misma masa total.
Esta redundancia es el gauge de la Sección~\ref{sec:predictor}, que aquí es la dirección
$(\mathbf1_R,-\mathbf1_C)$ de cada componente; lo fijamos omitiendo una restricción de columna por
componente.

Sea $p\in(0,1)^E$ con $Bp=\mathbf1$, donde $B$ contiene las restricciones RC independientes que
quedan tras esa omisión. Sean $D$ la incidencia casilla--diagonal, $h=\mathbf1-D^Tp>0$,
$q=\mathbf1-p$ y
\[
 H_B(p)=-\sum p\log p+3\sum q\log q-\sum h\log h.
\]
Introducimos actividades auxiliares $\lambda_c>0$, una por casilla, y escribimos $u=\log\lambda$,
componente a componente. Sea $p(u)$ una rama estacionaria de
$H_B(p)+u^Tp$ con hessiana definida negativa sobre $\ker B$.
Un máximo local estricto pero degenerado no basta.
Si $Z$ es una base ortonormal de $\ker B$, ponemos
\[
 a_c=1/p_c-3/q_c,\quad
 K=\operatorname{diag}a+D\operatorname{diag}(h^{-1})D^T,\quad
 K_T=Z^TKZ\succ0,\quad R=ZK_T^{-1}Z^T.
\]
El teorema de la función implícita da una rama diferenciable local, no un máximo global.
Al derivar $Z^T(\nabla H_B+u)=0$ se obtiene $-K_T\,dz+Z^Tdu=0$.
Por tanto $dp/du=R=R^T$ y $BR=0$.

Sea $d$ el número de componentes, constante mientras el soporte no cambia.
Usamos el determinante canónico del residual:
\[
 N=\begin{pmatrix}
 \operatorname{diag}a&D&B^T\\ D^T&-\operatorname{diag}h&0\\ B&0&0
 \end{pmatrix},\quad
 \mathfrak D=(-1)^{s+r}\prod_c(p_c/q_c^3)\det N>0,
\]
donde $s$ es el número de diagonales y $r=\operatorname{rank}B$.
Definimos $\ell=-\tfrac12\log\mathfrak D$, $g=\nabla_p\ell$ y
\[
 L(u)=H_B(p(u))+u^Tp(u)+\ell(p(u))+\tfrac d2\log(2\pi).
\]
El predictor del residual sin actividades auxiliares es $L(B)=L(0)$, con el mismo cofactor y la
misma normalización por componente.
El Lema~\ref{lem:puente} expresa $\mathfrak D$ mediante la matriz de líneas de la
Sección~\ref{sec:B3T05} y muestra que, en el tablero completo, $L(B)=\log Q_{\rm pred}(n)$.
\paragraph{Derivada total.}
La estacionariedad anula $(\nabla H_B+u)^Tdp$ porque $dp\in\ker B$.
Así,
\[
 \boxed{\nabla_uL=p+Rg},\qquad
 \boxed{\delta_c=(Rg)_c/p_c},\qquad
 \partial_{u_c}L=p_c(1+\delta_c).
\]
La derivada total del lazo es $p_c\delta_c$, no $\delta_c$.
La «marginal corregida» de la última igualdad es la derivada del predictor; no se identifica
con la marginal exacta del conteo.

\paragraph{Fórmula de Jacobi para el gradiente, sin polos en $p=1/4$.}
Sea $a'_c=-p_c^{-2}-3q_c^{-2}$ y pongamos los índices de las diagonales a continuación de los $E$
índices de casilla. Como $\partial_{p_c}(-h_\alpha)=D_{c\alpha}$,
\[
 \boxed{g_c=-\frac12\left[
 p_c^{-1}+3q_c^{-1}+a'_c(N^{-1})_{cc}
 +\sum_\alpha D_{c\alpha}(N^{-1})_{E+\alpha,E+\alpha}
 \right].}
\]
Esta expresión resulta de aplicar la fórmula de Jacobi a $\log|\det N|$, donde el determinante
tiene signo constante en la rama. En ella no se divide por $1-4p_c$. Además, $(N^{-1})_{pp}=R$
por el complemento de Schur diagonal y las restricciones RC.
Se puede obtener $Rg$ resolviendo $Nw=(g,0,0)^T$ y tomando su bloque inicial.
Una extensión del gradiente fuera de $Bp=1$ sólo añade componentes normales,
que $R$ anula; la identidad usada es intrínseca al espacio factible.

\paragraph{Normalización exacta.}
Bajo la traslación $u\mapsto u+B^T\eta$, el objetivo cambia en la constante
$\eta^TBp=\eta^T\mathbf1$. La rama correspondiente mantiene $p$, $h$,
el determinante y el número de componentes. Por tanto
\[
 L(u+B^T\eta)=L(u)+\eta^T\mathbf1,
 \quad B\nabla_uL=\mathbf1,
 \quad B(p\odot\delta)=BRg=0.
\]
La restricción eliminada en cada componente se reconstruye por la
igualdad de las sumas de filas y columnas. En toda fila original,
\[
 \boxed{\sum_{c\in\mathrm{fila}}p_c\delta_c=0.}
\]
Esta identidad no dice sólo que la suma sea $O(k^{-2})$, ni exige una aproximación asintótica.
Tampoco significa que $\sum_c\delta_c=0$.

\paragraph{Condiciones suficientes para una tasa uniforme.}
Para obtener $\|\delta\|_\infty=O(k^{-1})$ basta, por ejemplo,
\[
 \|\operatorname{diag}(p)^{-1}R\|_{\infty\to\infty}\le C_R,
 \qquad \inf_\eta\|g-B^T\eta\|_\infty\le C_g/k.
\]
Las Secciones~\ref{sec:B3T05} y~\ref{sec:B3T06} prueban estas cotas en la clase con brecha RC
ponderada uniforme.
